\needspace{12\baselineskip}
\SourceLeaf{313}
\begin{center}
\pdfbookmark[1]{11. Zu Art. 182–184.}{zusatz-011}
[11. Zu Art. 182--184.]
\end{center}
\begin{center}
\small
\rule{\textwidth}{.35pt}\par
Zeitschrift für Astronomie, herausgegeben von B. \textsc{von Lindenau} und J. G. F. \textsc{Bohnenberger}.\par
Erster Band. Heft für März und April 1816.\par
\rule{\textwidth}{.35pt}
\end{center}
Ich zeige bei dieser Gelegenheit noch eine Berichtigung an, die auf S. 218 und] 219 [S. 250--252 dieses Bandes] der Theoria Motus Corporum Coelestium vorzunehmen ist. Man lese nämlich
\[
\begin{aligned}
&\text{S. 218 Z. 3 statt}\quad e^{-\frac{hh\sigma\sigma}{\delta'''}}\quad\text{. . . .}\quad e^{-hh\delta'''\sigma\sigma}\\[5pt]
&\phantom{\text{S. 218 }}\text{Z. 4 statt}\quad\BareRoot{\frac{1}{\delta'''}}\quad\text{. . . . . .}\quad\BareRoot{\delta'''}\\[5pt]
&\text{S. 219 Z. 8 statt}\quad\BareRoot{A},\BareRoot{B'},\BareRoot{C''},\BareRoot{D'''}\\[4pt]
&\qquad\text{der Ordnung nach}\quad
\frac{1}{\BareRoot{A}}\mathpunct{\raisebox{.45ex}{,}}\quad
\frac{1}{\BareRoot{B'}}\mathpunct{\raisebox{.45ex}{,}}\quad
\frac{1}{\BareRoot{C''}}\mathpunct{\raisebox{.45ex}{,}}\quad
\frac{1}{\BareRoot{D'''}}\mathpunct{\cdot}
\end{aligned}
\]

Diese Unrichtigkeit hat ihren Ursprung darin, dass früher bei der Untersuchung andere Bezeichnungen verwendet worden waren, deren Austausch an den angegebenen Stellen versäumt wurde. Das numerische Beispiel ist, wie man sieht, gemäß den berichtigten Ausdrücken berechnet, doch wurde darin ein Rechenfehler begangen. Die Gleichung auf S. 219, Z. 5 von unten, muss nämlich lauten
\[
6633r=12707+2P-9Q+123R,
\]
folglich die Genauigkeit von \(r\) auf S. 220
\[
=\BareRoot{\frac{2211}{41}}=7{,}34.
\]

Herr \textsc{Nicolai} hat mich auf diese Berichtigungen aufmerksam gemacht; ich spreche ihm dafür hier meinen verbindlichsten Dank aus.
\begin{center}\rule{.18\textwidth}{.35pt}\end{center}
